\clearpage
\SourceLeaf{390}
\fancyhead[L]{\small MISCELLANEOUS PAPERS.}
\pdfbookmark[0]{Miscellaneous papers.}{verschiedene-aufsaetze}
\pdfbookmark[1]{Determination of the largest ellipse. 1810.}{groesste-ellipse-1810}
\begin{center}
{\Large MISCELLANEOUS PAPERS.}\par\bigskip
\rule{.85\textwidth}{.3pt}\par
v.\ \textsc{Zach}. Monatliche Correspondenz für Erd- und Himmelskunde. August 1810.\par
\rule{.85\textwidth}{.3pt}\par\bigskip
DETERMINATION OF THE LARGEST ELLIPSE\par\medskip
TANGENT TO THE FOUR SIDES OF A GIVEN QUADRILATERAL.\par\bigskip
\rule{.15\textwidth}{.3pt}
\end{center}
\hspace*{\parindent}I determine the positions of all points in the plane of the quadrilateral by abscissas and ordinates, first choosing the axis of abscissas and its origin entirely at will. I specify the quadrilateral by the points where its sides are met by the perpendiculars dropped to them from that origin, instead of specifying it by its vertices. Let these perpendiculars be \(a,a',a'',a'''\), and let their inclinations to the axis of abscissas be \(A,A',A'',A'''\), so that the coordinates of the four points of intersection just mentioned are
\[
\begin{array}{ll}
a\cos A,&a\sin A\\
a'\cos A',&a'\sin A'\\
a''\cos A'',&a''\sin A''\\
a'''\cos A''',&a'''\sin A'''
\end{array}
\]
